% GENERATED FILE. Edit canonical lessons, then run npm run generate:books.
% canonical-source-hash: fnv1a64:80769d3754463654
% canonical-lessons: ES-C641-taxi, ES-C641-taxista, ES-C641-tarifa, ES-C641-trayecto, ES-C641-parar

\chapter{En taxi}
\label{ch:es-en-taxi}
\hlchaptermodality{641}

\begin{chapteropening}
\textbf{By the end of this chapter:} \emph{I can take a taxi: ask about the fare, say how long the trip is and ask the driver to stop.}

Complete the last lesson of chapter 641: I can take a taxi: ask about the fare, say how long the trip is and ask the driver to stop.
\end{chapteropening}

\section[el taxi]{el taxi — the taxi}

\label{lesson:ES-C641-taxi}



\begin{quote}
Before the new one: say the Spanish for the van, then the Spanish for the driver.
\end{quote}

\subsection*{You'll want to know: el taxi}
\textbf{el taxi} — \textquotedblleft{}the taxi\textquotedblright{}.

\begin{quote}
Cogemos un taxi a la estación.
\end{quote}

\begin{quote}
\emph{We are taking a taxi to the station.}
\end{quote}

\emph{Coger un taxi} in Spain; in much of Latin America say \emph{tomar un taxi}, because \emph{coger} is vulgar there.

\begin{cousinweb}
\emph{taxi} is short for \emph{taxímetro}, which came from French \emph{taximètre}, a meter that measures the charge; it is built on Medieval Latin \emph{taxa} \textquotedblleft{}charge, rate\textquotedblright{} and Greek \emph{métron} \textquotedblleft{}measure\textquotedblright{} (EtymDB).
\end{cousinweb}

\begin{grammarlens}[title={What you've built}]
One more word for telling, explaining and describing at B1.
\end{grammarlens}

\subsection*{Your turn}
\begin{itemize}
\raggedright
  \item \emph{Say it:} \emph{el taxi}
  \item \emph{Say it:} \emph{el taxi}, once more
  \item \emph{Say it:} a sentence of your own with \emph{el taxi}
  \item \emph{Recall:} say the Spanish for the van, then the Spanish for the driver, then say \emph{el taxi} again
\end{itemize}

\begin{tcolorbox}[breakable,colback=teal!4,colframe=teal!35!black,title={Before you move on}]
What does el taxi mean? (\textquotedblleft{}The taxi\textquotedblright{}.) Say it once more.
\end{tcolorbox}
\section[el taxista]{el taxista — the taxi driver}

\label{lesson:ES-C641-taxista}



\begin{quote}
Before the new one: say the Spanish for the driver, then the Spanish for the taxi.
\end{quote}

\subsection*{You'll want to know: el taxista}
\textbf{el taxista} — \textquotedblleft{}the taxi driver\textquotedblright{}.

\begin{quote}
El taxista conoce todas las calles.
\end{quote}

\begin{quote}
\emph{The taxi driver knows every street.}
\end{quote}

One form for a man or a woman: \emph{el taxista}, \emph{la taxista}.

\begin{cousinweb}
\emph{taxista} is built on \emph{taxi} with \emph{-ista}, the ending for someone who does a job, which comes from Greek \emph{-ist\'{\={e}}s} through Latin.
\end{cousinweb}

\begin{grammarlens}[title={What you've built}]
One more word for telling, explaining and describing at B1.
\end{grammarlens}

\subsection*{Your turn}
\begin{itemize}
\raggedright
  \item \emph{Say it:} \emph{el taxista}
  \item \emph{Say it:} \emph{el taxista}, once more
  \item \emph{Say it:} a sentence of your own with \emph{el taxista}
  \item \emph{Recall:} say the Spanish for the driver, then the Spanish for the taxi, then say \emph{el taxista} again
\end{itemize}

\begin{tcolorbox}[breakable,colback=teal!4,colframe=teal!35!black,title={Before you move on}]
What does el taxista mean? (\textquotedblleft{}The taxi driver\textquotedblright{}.) Say it once more.
\end{tcolorbox}
\section[la tarifa]{la tarifa — the fare, the rate; the price list}

\label{lesson:ES-C641-tarifa}



\begin{quote}
Before the new one: say the Spanish for the taxi, then the Spanish for the taxi driver.
\end{quote}

\subsection*{You'll want to know: la tarifa}
\textbf{la tarifa} — \textquotedblleft{}the fare, the rate; the price list\textquotedblright{}.

\begin{quote}
La tarifa de noche es más cara.
\end{quote}

\begin{quote}
\emph{The night fare is more expensive.}
\end{quote}

For a phone or for electricity too: \emph{una tarifa plana} is a flat rate.

\begin{cousinweb}
\emph{tarifa} comes from Arabic \emph{taʿrīfa} \textquotedblleft{}a notice, a list of rates\textquotedblright{} (EtymDB), from the verb \emph{ʿarrafa} \textquotedblleft{}to make known\textquotedblright{}.
\end{cousinweb}

\begin{grammarlens}[title={What you've built}]
One more word for telling, explaining and describing at B1.
\end{grammarlens}

\subsection*{Your turn}
\begin{itemize}
\raggedright
  \item \emph{Say it:} \emph{la tarifa}
  \item \emph{Say it:} \emph{la tarifa}, once more
  \item \emph{Say it:} a sentence of your own with \emph{la tarifa}
  \item \emph{Recall:} say the Spanish for the taxi, then the Spanish for the taxi driver, then say \emph{la tarifa} again
\end{itemize}

\begin{tcolorbox}[breakable,colback=teal!4,colframe=teal!35!black,title={Before you move on}]
What does la tarifa mean? (\textquotedblleft{}The fare, the rate; the price list\textquotedblright{}.) Say it once more.
\end{tcolorbox}
\section[el trayecto]{el trayecto — the journey, the stretch (from one place to another)}

\label{lesson:ES-C641-trayecto}



\begin{quote}
Before the new one: say the Spanish for the taxi driver, then the Spanish for the fare.
\end{quote}

\subsection*{You'll want to know: el trayecto}
\textbf{el trayecto} — \textquotedblleft{}the journey, the stretch (from one place to another)\textquotedblright{}.

\begin{quote}
El trayecto dura veinte minutos.
\end{quote}

\begin{quote}
\emph{The journey takes twenty minutes.}
\end{quote}

One leg of the way; \emph{el viaje} is the whole trip.

\begin{cousinweb}
\emph{trayecto} comes from Latin \emph{trāiectus} \textquotedblleft{}a crossing over, a passage\textquotedblright{} (Lewis \& Short), from \emph{trāicere} \textquotedblleft{}to throw across\textquotedblright{}, built from \emph{trans} and \emph{iacere} \textquotedblleft{}to throw\textquotedblright{}.
\end{cousinweb}

\begin{grammarlens}[title={What you've built}]
One more word for telling, explaining and describing at B1.
\end{grammarlens}

\subsection*{Your turn}
\begin{itemize}
\raggedright
  \item \emph{Say it:} \emph{el trayecto}
  \item \emph{Say it:} \emph{el trayecto}, once more
  \item \emph{Say it:} a sentence of your own with \emph{el trayecto}
  \item \emph{Recall:} say the Spanish for the taxi driver, then the Spanish for the fare, then say \emph{el trayecto} again
\end{itemize}

\begin{tcolorbox}[breakable,colback=teal!4,colframe=teal!35!black,title={Before you move on}]
What does el trayecto mean? (\textquotedblleft{}The journey, the stretch (from one place to another)\textquotedblright{}.) Say it once more.
\end{tcolorbox}
\section[parar]{parar — to stop}

\label{lesson:ES-C641-parar}



\begin{quote}
Before the new one: say the Spanish for the fare, then the Spanish for the journey.
\end{quote}

\subsection*{You'll want to know: parar}
\textbf{parar} — \textquotedblleft{}to stop\textquotedblright{}.

\begin{quote}
¿Puede parar aquí, por favor?
\end{quote}

\begin{quote}
\emph{Can you stop here, please?}
\end{quote}

\emph{Pararse} is to come to a stop; in Latin America it also means to stand up.

\begin{cousinweb}
\emph{parar} is inherited from Latin \emph{parāre} \textquotedblleft{}to make or get ready, to prepare\textquotedblright{} (Lewis \& Short; EtymDB).
\end{cousinweb}

\begin{grammarlens}[title={What you've built}]
One more word for telling, explaining and describing at B1.
\end{grammarlens}

\subsection*{Your turn}
\begin{itemize}
\raggedright
  \item \emph{Say it:} \emph{parar}
  \item \emph{Say it:} \emph{parar}, once more
  \item \emph{Say it:} a sentence of your own with \emph{parar}
  \item \emph{Recall:} say the Spanish for the fare, then the Spanish for the journey, then say \emph{parar} again
\end{itemize}

\begin{tcolorbox}[breakable,colback=teal!4,colframe=teal!35!black,title={Before you move on}]
What does parar mean? (\textquotedblleft{}To stop\textquotedblright{}.) Say it once more.
\end{tcolorbox}
